% Generated draft: requires mathematical and source review.
\documentclass[11pt]{article}
\usepackage[margin=25mm]{geometry}
\usepackage{fontspec}
\setmainfont{Latin Modern Roman}
\usepackage{amsmath,amssymb,mathtools,unicode-math}
\setmathfont{Latin Modern Math}
\usepackage{enumitem,needspace,xurl,hyperref}
\setlength{\parindent}{0pt}
\setlength{\parskip}{5pt}
\setlength{\emergencystretch}{3em}
\newcommand{\N}{\mathbb{N}}
\newcommand{\Z}{\mathbb{Z}}
\newcommand{\Q}{\mathbb{Q}}
\newcommand{\R}{\mathbb{R}}
\newcommand{\C}{\mathbb{C}}
\newcommand{\F}{\mathbb{F}}
\newcommand{\A}{\mathbb{A}}
\newcommand{\PP}{\mathbb{P}}
\newcommand{\E}{\mathbb{E}}
\newcommand{\eps}{\varepsilon}
\newcommand{\dd}{\,\mathrm{d}}
\begin{document}
\section*{UnsolvedMath ID: 30001699}
\textbf{Source problem code:} catalogue entry
\subsection{Definitions and mathematical statement}
The supplied record is marked \texttt{unrecoverable}: its verification note says that it repeats source problem OWR-4798-005 and that no separate clean statement is promoted. The following is therefore a conditional formalization of the retained wording, not a recovered distinct problem.

Under the conventional interpretation of polytope diameter, let \(P\subseteq\R^d\) be a bounded, full-dimensional convex polytope. Its graph \(G(P)\) has the vertices of \(P\) as vertices and its edges as graph edges. Define
\[
\operatorname{diam}(P)=\max_{u,v\in V(P)}\operatorname{dist}_{G(P)}(u,v),
\]
where graph distance counts edges in a shortest path. For integers \(d\geq1\) and \(n\geq d+1\), define
\[
\Delta(d,n)=\max\{\operatorname{diam}(P):P\subseteq\R^d\text{ is a bounded, full-dimensional convex polytope with exactly }n\text{ facets}\}.
\]
These conventions concerning the ambient field, boundedness, graph diameter, and positive dimension are standard interpretations, but are not explicitly specified in the retained question.

The polynomial-bound question is whether
\[
\exists C>0\;\exists k\in\Z_{\geq0}\;\forall d\in\Z_{\geq1}\;\forall n\in\Z_{\geq d+1},\qquad
\Delta(d,n)\leq C(n+d)^k.
\]
Its superpolynomial alternative is the negation of this statement:
\[
\forall C>0\;\forall k\in\Z_{\geq0}\;\exists d\in\Z_{\geq1}\;\exists n\in\Z_{\geq d+1},\qquad
\Delta(d,n)>C(n+d)^k.
\]
The retained wording also asks whether the maximum diameter can be superlinear. If linear growth is measured against \(n+d\), this asks whether
\[
\forall C>0\;\exists d\in\Z_{\geq1}\;\exists n\in\Z_{\geq d+1},\qquad
\Delta(d,n)>C(n+d).
\]
The source record does not specify this meaning of superlinear or a particular regime in which \(n\) and \(d\) grow. Superlinear growth is compatible with a polynomial upper bound, so these are not mutually exclusive alternatives. The background's reference to ``simplex or polytope diameter'' does not specify an additional simplex-algorithm quantity.
\subsection{Short English statement}
Under the usual graph interpretation, ask whether the largest diameter of a convex \(d\)-polytope with \(n\) facets has a uniform polynomial bound in \(n\) and \(d\). The retained question also mentions superlinear growth without specifying its parameter regime. This record is a provenance-preserving duplicate marked \texttt{unrecoverable}; its supplied open label is not independently verified.
\subsection{Sources}
\begin{itemize}
\item Local UnsolvedMath registry, \texttt{data/unsolvedmath/problems.json}, permanent ID \(30001699\), exact source problem code \texttt{catalogue entry}, title ``Duplicate: Polynomial Bounds for Polytope Diameters.'' The supplied verification note identifies this as a repetition of \texttt{OWR-4798-005}; it states that no separate clean statement is promoted. No individual proposer is supplied.
\item \emph{Topological and Geometric Combinatorics} (2011), Oberwolfach Reports, pp.\ 349--423, DOI 10.4171/owr/2011/08. Supplied source URL and literature-source link, labeled ``original Oberwolfach report'': \url{https://doi.org/10.4171/owr/2011/08}. The supplied background describes Oberwolfach Reports as documenting mathematical workshops and their open-problem discussions. No problem-specific page or speaker attribution is supplied.
\item Literature assessment dated exactly \texttt{2026-08-21}: supplied status ``open.'' The supplied assessment says that the 2011 report leaves open whether simplex or polytope diameter can be superlinear or superpolynomial and that the cited lower bounds do not settle that question. This is supplied dated curation, not independently verified literature assessment or source verification. The supplied background explicitly requests independent verification of time-sensitive status and cited claims.
\end{itemize}
\end{document}
